\chapter{مقدمه}


در سال‌های اخیر، مدل‌های مولد\LTRfootnote{Generative Models} به‌عنوان یکی از ابزارهای قدرتمند در یادگیری ماشین و هوش مصنوعی، انقلابی در تولید داده‌های جدید ایجاد کرده‌اند. این مدل‌ها قابلیت تولید داده‌های متنوع و باکیفیت در حوزه‌های گوناگون مانند تصویر، متن، صوت و ویدئو را دارند و با الگوریتم‌های پیشرفته‌ای که به کار می‌گیرند، امکان بازسازی یا تولید داده‌هایی را فراهم می‌کنند که مورد نیاز افراد متخصص در طیف متنوعی از حوزه‌هاست. در میان این مدل‌ها، مدل‌های پخشی\LTRfootnote{Diffusion Models}  به‌عنوان یکی از جدیدترین و مؤثرترین خانواده‌های مدل‌های مولد حوزه‌ی تصویر معرفی شده‌اند که قابلیت‌های بی‌نظیری را در تولید داده‌ها نشان داده‌اند و به عنوان بهترین مدل\LTRfootnote{State of the Art} در حوزه‌ی تولید تصویر جایگاه خود را تثبیت کرده‌اند.
\\

مدل‌های پخشی با فرآیندی مبتنی بر تخریب تدریجی داده‌های واقعی و سپس بازسازی معکوس آن‌ها عمل می‌کنند. این مدل‌ها انواع مختلفی دارند از جمله مدل‌های پخشی مبتنی بر امتیاز\LTRfootnote{Score-based Diffusion Models}، مدل‌های پخشی احتمالاتی\LTRfootnote{Denoising Diffusion Probabilistic Models} و مد‌ل‌های پخشی مبتنی بر معادله دیفرانسیل\LTRfootnote{Stochastic Differential Equation based Models}. در مدل‌های پخشی احتمالاتی نویززدا، داده‌های واقعی به‌تدریج و در طول یک زنجیره‌ی مارکفی نویزی می‌شوند و در نهایت به یک توزیع ساده، معمولاً توزیع گاوسی، تبدیل می‌شوند. سپس در فرآیند رو به عقب\LTRfootnote{Reverse Process} با شروع از یک نویز خالص، نویززدایی داده‌ به‌صورت معکوس انجام می‌شود و به تولید نمونه‌های جدید منجر می‌گردد. این رویکرد ساده اما قدرتمند، مدل‌های پخشی را به یکی از برجسته‌ترین ابزارها برای تولید داده‌های باکیفیت تبدیل کرده است.  
\\

به طور کلی در هر مدل مولد حوزه‌ی تصویر سه نیازمندی مطرح است و اگر مدلی توانایی برآوردن این سه نیازمندی را داشته باشد می‌تواند در تمامی حوزه‌ها و تمامی کاربردها به راحتی مورد استفاده قرار گیرد. این سه نیازمندی عبارتند از کیفیت تصاویر تولیدی، تنوع تصاویر تولیدی و سرعت فرآیند تولید نمونه. مدل‌های پخشی در مقایسه با سایر مدل‌های مولد حوزه‌ی تصویر دارای کیفیت و تنوع نمونه‌های تولیدی هستند اما نمونه‌برداری در این مدل‌ها به نسبت سایر مدل‌ها بسیار آهسته است. فرآیند نمونه‌برداری در این مدل‌ها به گام‌های متعددی نیاز دارد و در هر گام باید محاسبات پیچیده‌ای انجام شود. این موضوع باعث می‌شود زمان اجرای مدل‌های پخشی برای تولید نمونه‌های جدید طولانی باشد و این امر استفاده از آن‌ها را در کاربردهای بلادرنگ و تعاملی محدود می‌کند.
\\

تحقیقات بسیاری در سال‌های اخیر برای غلبه بر این محدودیت و افزایش سرعت مدل‌های پخشی انجام شده است. این رویکردها را می‌توان به اشکال مختلف دسته‌بندی کرد اما یک دسته‌بندی رایج که در این پژوهش نیز مدنظر است تقسیم رویکردها بر اساس محل اثر آنهاست. رویکرد‌های افزایش سرعت می‌توانند با کاهش تعداد گام‌های مدل پخشی سرعت آن را افزایش دهند یا با کاهش حجم محاسباتی که در هر گام صورت می‌گیرد. رویکردهایی نظیر تقطیر دانش\LTRfootnote{Knowledge Distillation} و مدل‌های پخشی نویززدای ضمنی\LTRfootnote{Denoising Diffusion Implicit Models(DDIM)} در دسته‌ی کاهش تعداد گام‌ها قرار می‌گیرند. در مقابل، روش‌هایی مانند فشرده‌سازی شبکه‌ی نویززدا\LTRfootnote{Network Compression} و استفاده از حافظه‌ی پنهان\LTRfootnote{Caching} در دسته‌ی کاهش محاسبات هر گام قرار می‌گیرند. این روش‌ها اغلب با حفظ کیفیت خروجی، به دنبال کاهش زمان محاسبات و افزایش سرعت اجرای مدل بوده‌اند.  
\\

در این پژوهش با تمرکز بر روش‌های بدون آموزش کاهش محاسبات در هر گام، روش پیشنهادی تحت عنوان ادغام توکن‌های تطبیقی مبتنی بر حافظه‌ی پنهان ارائه شده است. این روش با تمرکز بر نقاط ضعف موجود در روش ادغام توکن‌ها\LTRfootnote{Token Merging} و همچنین با در نظر گرفتن این مشاهده که میزان تغییر بلوک‌های شبکه‌ی یوی نویززدا در گام‌های متوالی ناچیز و هموار است، تلاش می‌کند تا از طریق ترکیب روش‌ استفاده از حافظه‌ی پنهان و روش ادغام توکن‌ها میزان توکن‌های ورودی در هر بلوک مبدل\LTRfootnote{Transformer} را کاهش دهد. و این کاهش منجر به کاهش محاسبات انجام شده در هر گام از شبکه و در نتیجه کاهش زمان نمونه‌برداری در عین حفظ کیفیت می‌شود.
\\

نتایج آزمایشات انجام شده بر روی مجموعه داده‌ی $ImageNet$ نشان می‌دهد که روش پیشنهادی نسبت به روش پایه‌ی ادغام توکن‌ها و نیز نسبت به مدل پخشی $Stable\ Diffusion\ V1.5$ کاهش زمانی با ضریب ۲۴.۱ دارد و همچنین نشان داده‌ می‌شود که این کاهش حداقل کاهش زمانی‌ است که در روش ادغام توکن‌ها امکان دارد و این امر با ثابت نگه داشتن کیفیت تصاویر تولیدی انجام شده است.
\\

در ادامه این پژوهش، ابتدا مفاهیم پایه و مبانی نظری مرتبط با مدل‌های پخشی و روش‌های بهبود سرعت آن‌ها مورد بررسی قرار خواهد گرفت. سپس جزئیات روش پیشنهادی ارائه شده و مراحل پیاده‌سازی آن شرح داده می‌شود. در بخش‌های پایانی نتایج آزمایشات انجام‌شده تحلیل شده و کارایی روش پیشنهادی با سایر روش‌های موجود مقایسه می‌گردد. نهایتاً، راهکارهایی برای تحقیقات آینده در این زمینه ارائه خواهد شد.  
